\section{2D Vision 2}\label{d-vision-2}

\subsection{Object Detection}\label{object-detection}

\subsubsection{Definitions and Basic
Concepts}\label{definitions-and-basic-concepts}

\begin{itemize}
\tightlist
\item
  \textbf{Semantic
  Segmentation}：语义分割，逐像素预测类别，不区分同类不同实例。
\item
  \textbf{Instance
  Segmentation}：实例分割，逐像素预测类别，且区分同类不同实例。
\item
  \textbf{Object
  Detection}：介于二者之间，需要区分同类不同实例，但是不用精确到
  pixel，只需精确到 \textbf{tight bounding box}（bbox）。
\item
  \textbf{Bbox}：axis-aligned 的边界框，参数化为 \((x, y, h, w)\)，其中
  \(x, y\) 为左上角坐标。每个候选框还需附加一个 score 作为类别置信度。
\item
  \textbf{Backbone CNN}：特征提取网络，通常由 pretrained
  分类网络去除末层 FC 得到。输入原始图像，输出下采样后的 Feature Map。
\item
  \textbf{IoU}：定义边界框间的重叠程度，后续用于 RPN 样本划分、NMS
  及评估。 \[
    \text{IoU}(A,B) = \frac{\text{area}(A\cap B)}{\text{area}(A\cup B)}
    \] 其中 \(A,B\) 是两个边界框，\(\text{area}(\cdot)\)
  为框所覆盖的像素面积。
\end{itemize}

\subsubsection{Single Object Detection}\label{single-object-detection}

在 Backbone CNN 之后接两个 FC heads：

\begin{itemize}
\tightlist
\item
  \textbf{Classification Head}：输出 \(K\) 维向量，经 Softmax 得到 \(K\)
  个类别的概率分布，使用 \textbf{Cross-Entropy Loss}（Softmax
  Loss）监督；
\item
  \textbf{Regression Head}：输出 \(4\) 维向量，表示 bbox
  coordinates，使用 \textbf{L2 Loss}（MSE Loss）监督。
\end{itemize}

训练时，两个损失加权相加作为总损失，联合进行优化。

\subsubsection{Regression Losses}\label{regression-losses}

定位部分本质上是连续坐标预测，所以是一个回归问题，称为 Regression
Head。其中 Regression Loss 的选择尤为重要。

记 \(\Delta_i\) 为第 \(i\) 个回归分量，则：

\begin{itemize}
\item
  \textbf{L1 Loss}： \[
    \ell_{\text{L1}}(\Delta) = \sum_i |\Delta_i|
    \] 优点为对大误差不敏感，robust； 缺点为梯度恒定，收敛困难。
\item
  \textbf{L2 Loss}： \[
    \ell_{\text{L2}}(\Delta) = \sum_i \Delta_i^2
    \] 优点为收敛性更好； 缺点为梯度随误差线性增大，对大误差敏感。
\item
  \textbf{RMSE}（Root Mean Squared Error）： \[
    \ell_{\text{RMSE}}(\Delta) = \sqrt{\frac{1}{N}\sum_i \Delta_i^2}
    \] 缺点为误差接近 0 时使梯度不稳定，存在奇点。
\item
  \textbf{Smooth L1 loss}： \[
    \text{smooth}_{L1}(x) =
    \begin{cases}
    0.5x^2, & \text{if } |x| < 1\\
    |x| - 0.5, & \text{otherwise}
    \end{cases}
    \] 也即当误差小于阈值（通常取 1）时使用 L2，利于收敛；当误差大时使用
  L1，利于鲁棒。
\end{itemize}

\textbf{注}：smooth L1 就是 Huber Loss。

\subsubsection{Multi-Object Detection
Overview}\label{multi-object-detection-overview}

多目标检测要求输出\textbf{可变数量}的边界框。

\paragraph{Proposals vs RoIs}\label{proposals-vs-rois}

先生成一组可能包含物体的候选框 \textbf{Proposals}，数量较多（约为
2k），经筛选后得到 \textbf{RoIs}（Regions of Interest），后续只对 RoI
进行操作。

\begin{itemize}
\tightlist
\item
  \textbf{R-CNN / Fast R-CNN}：Proposals 由 Selective Search
  生成（无置信度分数），直接投影后进入 RoI Pool。
\item
  \textbf{Faster R-CNN}：Proposals 由 RPN 生成（带 \textbf{objectness
  score}），经 NMS 与正负采样筛选后得到 RoIs。
\end{itemize}

\paragraph{Proposal Generation}\label{proposal-generation}

\begin{itemize}
\tightlist
\item
  \textbf{Selective Search}（R-CNN / Fast R-CNN）：基于手工特征子识别。
\item
  \textbf{RPN}（Faster R-CNN）：在 Backbone 输出的 Feature Map
  使用可学习网络识别。
\end{itemize}

\paragraph{Sliding-Window}\label{sliding-window}

在图像上以不同大小的窗口滑动，对每种大小的每个窗口进行 Single Object。
显然计算量很大，效率很低。

\paragraph{Detection Paradigms}\label{detection-paradigms}

\begin{itemize}
\tightlist
\item
  \textbf{Two-Stage Detectors}：先提取 Region
  Proposal，再对每个区域进行分类与回归（精度高）。
\item
  \textbf{Single-Stage Detectors}：直接在 feature map
  上密集预测类别与边框，避免 proposal 生成阶段（速度快）。
\end{itemize}

\subsubsection{Two-Stage Detectors}\label{two-stage-detectors}

核心思想：先提取 Region Proposal，再对每个 proposal
进行\textbf{类别分类}与\textbf{边框回归}。

\paragraph{R-CNN}\label{r-cnn}

Region-based CNN。

\begin{enumerate}
\def\labelenumi{\arabic{enumi}.}
\tightlist
\item
  用 Selective Search 生成 \textbf{Proposals}。
\item
  每个 \textbf{Proposal} 使用 bilinear 插值放缩到固定大小，然后带入
  Backbone CNN 中提取特征。
\item
  最后经过 cls head 和 reg head。
\end{enumerate}

\textbf{问题}：

\begin{itemize}
\tightlist
\item
  每个 Proposal 都要独立前向通过 Backbone CNN，计算重复且慢。
\item
  多个部件（CNN、Cls、Reg）分别训练，非 end-to-end。
\item
  cropped 的 Proposal 缺失全图 context，导致边框回归时
  \textbf{Information Inefficient}！
\end{itemize}

\paragraph{Fast R-CNN}\label{fast-r-cnn}

既然先 crop 后 conv 会丢失 context，不妨\textbf{先 conv 再 crop}。

\begin{enumerate}
\def\labelenumi{\arabic{enumi}.}
\tightlist
\item
  Backbone CNN 提取 feature map。
\item
  用 Selective Search 生成 \textbf{Proposals}。
\item
  将 \textbf{Proposals} 投影到 feature map
  上，边界吸附到最近的网格点中，得到对应的 \textbf{RoI}。
\item
  使用 \textbf{RoI Pool} 进行 resize。也即将这些 RoI 划分为固定数量的
  bin，在每个 bin 上做 max pooling。
\item
  最后经过 cls head 和 reg head。
\end{enumerate}

\begin{itemize}
\tightlist
\item
  优点：减少重复计算，显著加速；部分端到端，训练更方便。
\item
  缺点：速度受 Selective Search 制约；有一定精度误差。
\end{itemize}

\paragraph{Faster R-CNN}\label{faster-r-cnn}

核心在于采用 RPN 替代了 Selective Search，直接在共享的 feature map
上生成 proposals，从而实现近似的端到端训练。

\textbf{Pipeline}：

\begin{enumerate}
\def\labelenumi{\arabic{enumi}.}
\tightlist
\item
  Backbone CNN 提取 feature map。
\item
  用 RPN 在 feature map 上生成高质量 \textbf{Proposals}。
\item
  经 RoI Generation 筛选后的 \textbf{RoIs} 通过 RoI Pooling resize
  到固定尺寸。
\item
  最后经过 cls head 和 reg head。
\end{enumerate}

\begin{itemize}
\tightlist
\item
  训练：RPN 与 heads 共用 Backbone，可以交替训练或联合训练。
\item
  优点：速度快、候选更准确且可学习、整体性能提升显著。
\end{itemize}

\textbf{RoI Generation}：

RPN 输出的 Proposals 需经后处理才能成为 RoIs：

\begin{enumerate}
\def\labelenumi{\arabic{enumi}.}
\tightlist
\item
  \textbf{NMS}：按 RPN 预测的 \textbf{objectness score}
  降序，去除与高分框 IoU 过高的冗余框（通常 \(\tau=0.7\)）；
\item
  \textbf{Top-K}：保留分数最高的 \(K\) 个；
\item
  \textbf{正负采样}（train only）：计算与 GT 的 IoU，划分为 Positive 和
  Negative。按比例采样得到恒定数量的 RoIs，送入后续 Heads。
\end{enumerate}

\textbf{Faster R-CNN Paradigm}：

\begin{itemize}
\tightlist
\item
  \textbf{Run once per image}：Backbone 提取 feature map，RPN 生成
  Proposals 并筛选出 RoIs。
\item
  \textbf{Run once per region}：RoI Pool / Align，two-heads。
\end{itemize}

\paragraph{RPN}\label{rpn}

预定义一组固定尺寸与纵横比的参考框，称为 \textbf{Anchor}。 在 feature
map 的每个位置上放置多个 Anchors，以覆盖不同大小与形状的目标。

对每个 anchor，两个 heads 输出：

\begin{itemize}
\tightlist
\item
  \textbf{cls}：预测该 anchor 是否为前景，得到 \textbf{objectness
  score}。
\item
  \textbf{reg}：预测该 anchor 相对于真实框的\textbf{坐标偏移量}。
\end{itemize}

\textbf{Parameterization}：

由于尺寸、纵横比差异大，故需对回归的目标进行归一化改进，减小尺度差异带来的数值问题。

给定候选框 \((x_a, y_a, w_a, h_a)\) 和真实框
\((x_g, y_g, w_g, h_g)\)，回归目标通常定义为： \[
t_x = \frac{x_g - x_a}{w_a},\quad
t_y = \frac{y_g - y_a}{h_a},\quad
t_w = \log\frac{w_g}{w_a},\quad
t_h = \log\frac{h_g}{h_a}
\]

\textbf{Training}：

计算每个 anchor 与所有 GT 框的 IoU，按阈值划分样本：

\begin{itemize}
\tightlist
\item
  \textbf{Positive}：\(\text{IoU} > 0.7\)，或与某 GT 的 IoU 最大；
\item
  \textbf{Negative}：\(\text{IoU} < 0.3\)；
\item
  \textbf{Ignore}：\(\text{IoU} \in [0.3, 0.7]\)，不参与训练。
\end{itemize}

\textbf{Loss Function}：

\begin{enumerate}
\def\labelenumi{\arabic{enumi}.}
\item
  \textbf{Classification Loss}：

  对每个 anchor \(i\)，objectness confidence 为 \(p_i\)，one-hot label
  为 \(p_i^* \in \{0, 1\}\)。采用二分类交叉熵： \[
   L_{\text{cls}} = -\frac{1}{N_{\text{cls}}} \sum_i \Bigl[ p_i^* \log(p_i) + (1-p_i^*) \log(1-p_i) \Bigr]
   \]

  其中 \(N_{\text{cls}}\) 为参与分类训练的 anchor
  数量（通常为正负样本采样后的总数）。
\item
  \textbf{Regression Loss}：

  仅对 positive（\(p_i^* = 1\)）计算边框偏移。设预测偏移量为
  \(t_i = (t_x, t_y, t_w, t_h)\)，真实目标偏移为 \(t_i^*\)，则：

  \[
   L_{\text{reg}} = \frac{1}{N_{\text{reg}}} \sum_i p_i^* \cdot \text{Smooth}_{L1}(t_i - t_i^*)
   \]

  其中 \(\text{Smooth}_{L1}\) 定义为：

  \[
   \text{Smooth}_{L1}(x) =
   \begin{cases}
   0.5 x^2, & |x| < 1 \\[6pt]
   |x| - 0.5, & |x| \geq 1
   \end{cases}
   \] 负样本时 \(p_i^*=0\)，不参与回归。
\item
  \textbf{Total Loss}：

  将两部分加权求和：

  \[
   L(\{p_i\}, \{t_i\}) = L_{\text{cls}} + \lambda L_{\text{reg}}
   \]

  其中 \(\lambda\) 为平衡系数。
\end{enumerate}

\subsubsection{Single-Stage Detectors}\label{single-stage-detectors}

Single-Stage Detectors 直接在 feature map
上密集预测类别与边框，避免候选生成阶段，因而推理速度快。

代表方法包括 \textbf{YOLO}（You Only Look Once）、\textbf{SSD}（Single
Shot MultiBox Detector）、\textbf{RetinaNet} 等。

\textbf{Core Mechanism}：

将图像（或 feature map）划分为 \(H \times W\) 的网格。在每个 grid cell
上预置 \(B\) 个不同尺度与纵横比的 \textbf{anchors}。网络直接输出：

\begin{itemize}
\tightlist
\item
  \textbf{reg}：对每个 anchor 预测中心偏移
  \((\Delta x, \Delta y)\)、尺寸偏移
  \((\Delta h, \Delta w)\)，以及该框的 objectness。
\item
  \textbf{cls}：对 \(K\) 个类别（包含背景）输出 score。
\end{itemize}

输出张量维度为 \(H \times W \times (5B + C)\)。

本质上，RPN 只区分前景和背景，而 Single-Stage Detector 是
\textbf{category-specific}。

\subsubsection{NMS}\label{nms}

使用 NMS 删除重复且重叠的检测框，确保 single response constraint。

\textbf{Input}：候选框集合 \(B\)，对应的 class score \(S\)，IoU 阈值
\(\tau\)。\\
\textbf{Output}：最终检测结果集合 \(D\)。

\textbf{Algorithm}：

\begin{enumerate}
\def\labelenumi{\arabic{enumi}.}
\tightlist
\item
  初始化：\(D \leftarrow \varnothing\)。
\item
  按 class score \(S\) 将 \(B\) 中所有候选框\textbf{降序排序}。
\item
  从 \(B\) 中选取当前分数最高的框 \(b\)，将其从 \(B\) 移除，并加入
  \(D\)。
\item
  将 \(b\) 与 \(B\) 中剩余所有候选框逐一计算 IoU，若
  \(\text{IoU}(b, b') > \tau\)，则将 \(b'\) 从 \(B\) 中移除。
\item
  重复步骤 3--4，直到 \(B\) 为空。
\item
  返回 \(D\)。
\end{enumerate}

\textbf{Variants}：

\begin{itemize}
\tightlist
\item
  \textbf{Soft-NMS}：不直接删除高 IoU
  框，而是根据重叠程度衰减其置信度分数，从而降低硬删除造成的召回损失。
\item
  \textbf{Class-wise NMS}：对每个类别独立执行
  NMS，避免不同类别的检测框互相抑制。
\end{itemize}

\textbf{Notes}：

\begin{itemize}
\tightlist
\item
  Soft-NMS
  可以缓解多个真实目标在空间中确实高度重叠的问题，如拥挤、遮挡场景。
\item
  Class-wise NMS
  可以缓解不同类别的框即使空间重叠也应分别保留的问题，如人和背包的重叠。
\end{itemize}

\subsubsection{Evaluation Metrics}\label{evaluation-metrics}

检测任务的评估不仅要考虑 classification
accuracy，还要考虑定位精度和重复检测。核心指标为 \textbf{AP}（Average
Precision）与 \textbf{mAP}（mean AP）。

\textbf{Basic Definitions}：

\begin{itemize}
\tightlist
\item
  \textbf{Precision}： \[
    \text{Precision} = \frac{\text{TP}}{\text{TP} + \text{FP}}
    \]
\item
  \textbf{Recall}： \[
    \text{Recall} = \frac{\text{TP}}{\text{TP} + \text{FN}}
    \]
\item
  \textbf{TP}：类别正确，\(\text{IoU}(\text{pred}, \text{gt}) \ge \tau\)，且该
  GT 框此前未被其他更高置信度的预测框匹配过。
\item
  \textbf{FP}：类别错误，\(\text{IoU}(\text{pred}, \text{gt}) < \tau\)，或匹配到已被占用的
  GT 的预测框。
\item
  \textbf{FN}：未被预测框匹配的 GT 框。
\end{itemize}

\textbf{AP}（Average Precision）：

对单个类别，将检测结果按 class score
降序，计算不同阈值下（依次遍历每个检测结果的 class score 即可）的
precision 与 recall，绘制 Precision-Recall 曲线。

AP 是 PR 曲线下的面积。通常采用 11-point 插值来算： \[
\text{AP} = \frac{1}{11}\sum_{\text{Recall}_i} \text{Precision}(\text{Recall}_i)
\]

\textbf{mAP}（mean AP）：

对多个类别取平均，或对不同 IoU 阈值取平均，或对类别和阈值一起取平均。

\begin{center}\rule{0.5\linewidth}{0.5pt}\end{center}

\subsection{Instance Segmentation: Mask
R-CNN}\label{instance-segmentation-mask-r-cnn}

实例分割，逐像素预测类别，且区分同类不同实例。

主流思路分为 bottom-up 与 top-down 两类。

方法 \textbar{} 思路 \textbar{} 优点 \textbar{} 缺点 \textbar{}\\
:--- \textbar{} :--- \textbar{} :--- \textbar{} :--- \textbar{}\\
\textbf{Bottom-up} \textbar{} 先做 pixel grouping，再将每个 group
标注为不同实例 \textbar{} 不依赖 object detector，能处理重叠严重物体
\textbar{} pixel grouping 困难 \textbar{}\\
\textbf{Top-down} \textbar{} 先做 object detection，再在每个框内标记
binary mask \textbar{} 与 object detector 结合紧密，直观可扩展
\textbar{} 依赖 object detection 质量 \textbar{}\\
\textbf{Hybrid} / \textbf{Proposal-free} \textbar{} 直接做
transformer-based 的集合预测 \textbar{} 可实现 end-to-end \textbar{}
设计与训练较复杂 \textbar{}

\textbf{Mask R-CNN} 是一种典型的 \textbf{Top-down} 框架，基于 Faster
R-CNN 结构，额外增加了一个 mask head 用于预测每个 RoI 的 binary mask。

\subsubsection{Pipeline}\label{pipeline}

\begin{enumerate}
\def\labelenumi{\arabic{enumi}.}
\tightlist
\item
  Backbone CNN 提取 feature map。
\item
  使用 RPN 在 feature map 上生成 proposals。
\item
  对每个 proposal 使用 RoIAlign 从 feature map 中裁剪固定大小的 RoI。
\item
  将 RoI 特征分别送入三个并行分支：

  \begin{itemize}
  \tightlist
  \item
    cls head：预测 classification。
  \item
    reg head：回归 bbox。
  \item
    mask head：为该 RoI 预测一个同分辨率的 binary mask。
  \end{itemize}
\end{enumerate}

\textbf{注}：在训练时，只对 positive 进行 mask head。

\subsubsection{RoIAlign}\label{roialign}

原始的 RoI Pool 在将 proposal 区域映射到 feature map 并划分 bin
时，会进行离散化（quantization），导致误差。

为解决这一问题，Mask R-CNN 引入了 \textbf{RoIAlign}：

\begin{enumerate}
\def\labelenumi{\arabic{enumi}.}
\tightlist
\item
  Proposal project 得到的浮点位置不变，直接划分 bin。
\item
  对每个 bin 采样若干点（常用 4 个点），对采样点使用 bilinear
  interpolation 进行取值，然后做 max 或 average pooling。
\item
  输出固定大小的 RoI。
\end{enumerate}

\textbf{注}：这一看似微小的改进能大大提升精度，尤其是高阈值时的 AP
提升明显。

\subsubsection{Mask Head Design}\label{mask-head-design}

\paragraph{Output: Multinomial vs
Independent}\label{output-multinomial-vs-independent}

\textbf{Multinomial}（per pixel Softmax）：

标准语义分割的做法，对每个像素输出 \(K\) 维向量并经 Softmax，得到 \(K\)
个类别的概率分布。 但在 Mask R-CNN 中，类别已经\textbf{由 cls head
单独判定}，mask head 只需判断前景/背景。 若使用
Softmax，不同类别的掩码概率会相互挤压（class
competition），导致掩码预测与分类结果冲突。

\textbf{Independent}（per pixel Sigmoid）：Mask R-CNN
采用此方式。每像素独立做二分类，与类别完全解耦，避免概率竞争，使 mask
head 专注于实例形状。

\paragraph{Architecture: Class-specific vs
Class-agnostic}\label{architecture-class-specific-vs-class-agnostic}

在确定了 Sigmoid
二分类的形式后，还需决定是否为每个类别维护独立的掩码通道。

\textbf{Class-specific}：输出 \(K\) 通道，也即为 \(K\) 个类别各预测一个
binary mask。

\begin{itemize}
\tightlist
\item
  优点：自然编码类别相关的形状先验（如人与车的形状不同）。
\item
  缺点：参数量随类别数线性增长，类别较多时代价高。
\end{itemize}

\textbf{Class-agnostic}：输出 \(1\) 通道，也即所有类别共享同一个 binary
mask。

\begin{itemize}
\tightlist
\item
  优点：参数量小，训练更简洁。在许多场景下性能接近 class-specific。
\end{itemize}

Mask R-CNN 采用 \textbf{class-specific} 的 \(K\)
通道输出，但训练与推理时仅使用 cls head 预测的类别所对应的那个通道。

\paragraph{Loss: Per-Pixel Binary CE}\label{loss-per-pixel-binary-ce}

Mask R-CNN 对每个 RoI 的掩码使用 \textbf{Per-Pixel Binary
Cross-Entropy}： \[
\mathcal{L}_{mask}
= -\frac{1}{m^2} \sum_{u=1}^{m} \sum_{v=1}^{m}
\left[ y_{uv} \log \hat{y}_{uv} + (1 - y_{uv}) \log (1 - \hat{y}_{uv}) \right]
\]

其中：

\begin{itemize}
\tightlist
\item
  \(m\)：掩码分辨率（如 \(28 \times 28\)），与 RoI 的分辨率相同；
\item
  \(y_{uv}\in\{0,1\}\)：对应位置的真实掩码像素；
\item
  \(\hat{y}_{uv}\in(0,1)\)：网络预测的该位置属于该实例的概率；
\end{itemize}

\textbf{注}：训练时\textbf{仅对该 RoI 的 GT 类别所对应的通道计算 BCE
Loss}，其余通道不参与梯度回传，避免了不同类别掩码之间的梯度干扰。
